\section{Derivation of the Gamma product for the residue}
\label{app:product}

The following determinant is a classical finite algebraic ingredient of
the source's Gauss--Vandermonde proof. We include its derivation so the
normalization in \eqref{eq:residue-product} can be checked independently.

\begin{lemma}[A rational factorial determinant]\label{lem:factorial-det}
For indeterminates $x_0,\ldots,x_{L-1},t$,
\begin{equation}\label{eq:factorial-det}
 \det_{0\le i,j<L}\frac{(x_i)_j}{(x_i+t)_j}
 =\frac{\displaystyle\prod_{j=0}^{L-1}(t)_j
             \prod_{0\le i<k<L}(x_k-x_i)}
        {\displaystyle\prod_{i=0}^{L-1}(x_i+t)_{L-1}}.
\end{equation}
It is an identity of rational functions and hence holds wherever the
specialized expressions are defined, with removable values interpreted
by cancellation.
\end{lemma}
\begin{proof}
Multiply row $i$ by $(x_i+t)_{L-1}$. Column $j$ becomes the polynomial
\[
 p_j(x)=(x)_j(x+t+j)_{L-1-j},
\]
of degree at most $L-1$. Its evaluation determinant is the Vandermonde
in the $x_i$ times the determinant of the polynomial coefficients.
Evaluate at $x_i=-i$. The matrix is lower triangular, since $(-i)_j=0$
for $j>i$, and its diagonal is
$(-1)^ii!(t)_{L-1-i}$. Dividing its determinant by the Vandermonde
at $0,-1,\ldots,-(L-1)$ leaves $\prod_{j=0}^{L-1}(t)_j$.
Dividing out the original row factors proves the identity.
\end{proof}

Set $L=r-1$ and $d=2r-2$. Since the two polynomial bases
$\fall ni$ and $(n+a)_j$ are monic in their respective degrees,
Cauchy--Binet gives
\begin{equation}\label{eq:residue-mixed-moment}
 V_r(1,a)=\det_{0\le i,j<L}
 \sum_{n\ge0}\fall ni(n+a)_j
             \frac{\Gamma(n+a)}{\Gamma(n+a+d+1)}.
\end{equation}
Terms $n<i$ vanish. Shift $n=i+k$ and apply Gauss summation
\cite{dlmf:gauss}. Because $i+j\le2r-4=d-2$, its excess
$d-i-j$ is positive. The resulting entry is
\begin{equation}\label{eq:residue-entry}
 M_{ij}=\frac{\Gamma(i+1)\Gamma(a+i+j)\Gamma(d-i-j)}
                    {\Gamma(d+1-j)\Gamma(a+d)}.
\end{equation}
Factor the row terms $\Gamma(i+1)\Gamma(a+i)\Gamma(d-i)$,
the column terms $(-1)^j/\Gamma(d+1-j)$, and one factor
$1/\Gamma(a+d)$ per row. The remaining determinant is
\[
 \det_{0\le i,j<L}\frac{(a+i)_j}{(1-d+i)_j}.
\]
Apply Lemma~\ref{lem:factorial-det} with
$x_i=a+i$ and $t=1-d-a$. None of the denominator rising factorials
vanishes in the required range of indices. Its Vandermonde is
$\prod_{k=0}^{L-1}k!$. Substitution gives the intermediate product
\begin{align}
 V_r(1,a)={}&
 \frac{\prod_{i=0}^{L-1}\Gamma(i+1)\Gamma(a+i)\Gamma(d-i)}
      {\Gamma(a+d)^L}
 \prod_{j=0}^{L-1}\frac{(-1)^j}{\Gamma(d+1-j)}\notag\\
 &\times
 \frac{\prod_{j=0}^{L-1}(1-d-a)_j\prod_{k=0}^{L-1}k!}
      {\prod_{i=0}^{L-1}(1-d+i)_{L-1}}.
 \label{eq:residue-intermediate}
\end{align}
Using
\[
 (1-d-a)_j=(-1)^j\frac{\Gamma(a+d)}{\Gamma(a+d-j)},\qquad
 (1-d+i)_{L-1}=(-1)^{L-1}
       \frac{\Gamma(d-i)}{\Gamma(d-i-L+1)},
\]
and $d-L+1=r$, the signs cancel and
\eqref{eq:residue-intermediate} becomes
\[
 \prod_{j=0}^{r-2}
 \frac{j!\,\Gamma(a+j)\Gamma(j+1)\Gamma(r-j)}
       {\Gamma(2r-1-j)\Gamma(a+2r-2-j)}.
\]
This proves \eqref{eq:residue-product} for $\Rea a>0$. In particular,
\[
 V_2(1,a)=\frac1{2a(a+1)},\qquad
 V_3(1,a)=\frac1{72a(a+1)^2(a+2)^2(a+3)}.
\]
No additional $L!$ is introduced, since the defining ensemble uses
strictly ordered indices throughout.
